\chapter{Gender Incongruence}
\index{Gender Incongruence}

\section*{Definition and Overview}
Gender incongruence is the term used when a person's gender identity, their deeply felt sense of being male, female, or another gender, differs from the sex they were assigned at birth. It is a normal and valid human variation, not an illness or disorder, and the World Health Organization moved it out of the mental disorders chapter of its classification to reflect this. When gender incongruence causes significant distress, it is called gender dysphoria, and it is that distress, not the identity itself, that may benefit from care. Care is affirming, individualised, and may include psychological support, social affirmation, hormone therapy, and, for some people, surgery. This chapter explains gender incongruence, the experiences of gender-diverse people, and the care available.

\section*{Epidemiology}
Around 1--2% of the population is estimated to be transgender or gender diverse, and a larger number experiences some degree of gender incongruence. The proportion of young people identifying as transgender or non-binary has risen in recent years, with improved awareness and acceptance. Many gender-diverse people experience significant mental health disparities, including higher rates of anxiety, depression, and suicide, which are largely driven by stigma, discrimination, and barriers to care rather than by gender identity itself. Supportive families, schools, and health services markedly improve outcomes.

\section*{Aetiology and Pathophysiology}
\begin{itemize}
    \item \textbf{Gender identity:} a person's internal sense of their gender, which develops early in life and is influenced by biological, hormonal, and social factors
    \item \textbf{Gender incongruence:} the mismatch between identity and the sex recorded at birth
    \item \textbf{Non-binary identities:} some people do not identify as exclusively male or female, and identities span a spectrum
    \item \textbf{Gender dysphoria:} the clinically significant distress that incongruence can cause, including anxiety, depression, and discomfort with physical characteristics
    \item \textbf{Not a mental illness:} international classifications recognise that the identity itself is not a disorder, though the distress is real and treatable
    \item \textbf{Early development:} many transgender children identify their gender early, while others recognise it in adolescence or adulthood
    \item \textbf{The minority stress model:} stigma, rejection, and discrimination, not identity, drive much of the mental health burden
    \item \textbf{Affirming care:} supporting the person's identity reduces distress and improves wellbeing
\end{itemize}

\section*{Clinical Features}
\begin{itemize}
    \item A persistent sense that one's gender differs from the sex assigned at birth
    \item Distress about body characteristics, social roles, or being referred to by the wrong name or pronouns
    \item In children: insistence on being treated as the other gender, and distress about puberty
    \item Anxiety, depression, and withdrawal, often linked to rejection or concealment
    \item A desire for social affirmation, different name and pronouns, or medical treatment
    \item Many gender-diverse people live happy, fulfilling lives without any medical treatment
    \item The features vary enormously; there is no single way to experience gender incongruence
\end{itemize}

\section*{Differential Diagnosis}
Assessment distinguishes gender identity from other issues.
\begin{itemize}
    \item Body dysmorphic disorder, in which a person dislikes a specific feature without gender incongruence
    \item Autogynephilia and paraphilic concepts, which are outdated and not supported
    \item Mental health conditions, including depression, anxiety, and trauma, which may coexist and need treatment
    \item Autistic spectrum traits, which are overrepresented in gender-diverse people
    \item The key distinction is whether the distress arises from the gender identity itself or from other factors
    \item A respectful, unhurried assessment by an experienced professional resolves these
\end{itemize}

\section*{When to Seek Medical Care}
Gender-diverse people should have access to affirming health care. A GP can provide support, refer to specialist gender services, prescribe hormone therapy where appropriate, and coordinate mental health care. Families of gender-diverse children should seek guidance from affirming professionals.

\textbf{Contact your local emergency services for:}
\begin{itemize}
    \item Suicidal thoughts or behaviour, which are more common in unsupported gender-diverse people
    \item Severe distress, self-harm, or a crisis situation
    \item Any medical emergency
\end{itemize}

\textbf{Support lines:} QLife (1800 184 527) provides counselling and referral for LGBTQIA+ people, and a local crisis service is available 24 hours.

\section*{Diagnosis and Investigations}
\begin{itemize}
    \item \textbf{Assessment:} a supportive conversation about identity, distress, and needs, conducted without assumptions
    \item \textbf{Diagnosis of gender dysphoria:} made when the incongruence causes clinically significant distress, using the diagnostic criteria
    \item \textbf{Mental health assessment:} for coexisting depression, anxiety, and trauma
    \item \textbf{Family assessment:} in children and adolescents, involving the family supportively
    \item \textbf{Medical assessment:} for people considering hormone therapy, including blood tests
    \item \textbf{No testing:} gender identity is not diagnosed by tests; the person is the authority on their own identity
\end{itemize}

\section*{Management}
\begin{itemize}
    \item \textbf{Psychological support:} to explore identity, manage distress, and address stigma and discrimination
    \item \textbf{Social affirmation:} using the person's name and pronouns, and supporting their social transition
    \item \textbf{Family and school support:} education and advocacy to create accepting environments
    \item \textbf{Hormone therapy:} puberty blockers and hormone therapy for adolescents under specialist care, and hormone therapy for adults, which has well-documented benefits
    \item \textbf{Surgery:} available to adults who choose it, after assessment
    \item \textbf{Treat coexisting conditions:} depression, anxiety, and other issues are treated as for anyone
    \item \textbf{Peer support:} connection with LGBTQIA+ communities and services
    \item \textbf{No treatment needed:} many people need only support and acceptance, not medical treatment
\end{itemize}

\section*{Complications}
\begin{itemize}
    \item The mental health burden of stigma, rejection, and family conflict
    \item Delays in care, which worsen distress
    \item The harms of conversion practices, which are harmful and banned in many countries
    \item Barriers to employment, housing, and health care
    \item Suicidality in unsupported people
    \item The risks of medical treatments, which are discussed and monitored
\end{itemize}

\section*{Prognosis and Outcomes}
With affirming support, gender-diverse people live happy, healthy lives with outcomes no different from the general population. Affirming care is associated with reduced depression and anxiety and markedly lower suicidality, while rejection and conversion practices cause serious harm. Hormone therapy and surgery, for those who choose them, have high satisfaction. The strongest predictors of good outcomes are supportive families, affirming health care, and acceptance in schools and communities.

\section*{Prevention and Self-Care}
\begin{itemize}
    \item Create and seek accepting environments; support makes the difference
    \item Connect with affirming health services and LGBTQIA+ community organisations
    \item Use affirming names and pronouns for gender-diverse people
    \item Families should support their children's exploration and seek affirming professional advice
    \item Schools and workplaces should adopt inclusive policies
    \item Access mental health support early when distress arises
    \item If you are gender diverse, know your rights and where to find safe care
    \item Challenge stigma and discrimination in your community
\end{itemize}

\section*{Sources and Further Reading}
The following reputable sources provide further information for healthcare students and patients seeking reliable, current advice about gender incongruence.
\begin{itemize}
    \item TransHub (ACON). \textit{Gender identity information.} https://www.transhub.org.au/
    \item QLife. \textit{Counselling and support for LGBTQIA+ people.} https://qlife.org.au/
    \item AusPATH. \textit{Guidelines for gender-affirming care.} https://auspath.org.au/
    \item Healthdirect Australia. \textit{Gender dysphoria.} https://www.healthdirect.gov.au/
    \item World Professional Association for Transgender Health (WPATH). \textit{Standards of Care.} https://www.wpath.org/
    \item World Health Organization. \textit{Gender and health.} https://www.who.int/
\end{itemize}
